A pendulum of length $\ell$ is released from the horizontal position. It can be shown that the tangential $F = ma$ equation is (where $\theta$ is measured with respect to the vertical)

$$\ddot{\theta} = -\frac{g \sin \theta}{\ell}. \tag{1.20}$$

If $\ell = 1 \text{ m}$, and $g = 9.8 \text{ m/s}^2$, write a program to show that the time it takes the pendulum to swing down through the vertical position is $t \approx 0.592 \text{ s}$. This happens to be about $1.18$ times the $(\pi/2)\sqrt{\ell/g} \approx 0.502 \text{ s}$ it would take the pendulum to swing down if it were released from very close to the vertical (this is $1/4$ of the standard period of $2\pi\sqrt{\ell/g}$ for a pendulum). It also happens to be about $1.31$ times the $\sqrt{2\ell/g} \approx 0.452 \text{ s}$ it would take a mass to simply freefall a height $\ell$.